\documentclass{article}
\usepackage{amsmath,amssymb,amsthm}
\usepackage{algorithm}
\usepackage[noend]{algpseudocode}
\usepackage{geometry}
\geometry{margin=1in}
\newtheorem{theorem}{Theorem}
\newtheorem{lemma}{Lemma}
\newtheorem{assumption}{Assumption}
\newcommand{\R}{\mathbb{R}}
\newcommand{\C}{\mathcal{C}}
\newcommand{\gap}{\operatorname{gap}}
\newcommand{\dist}{\operatorname{dist}}
\newcommand{\conv}{\operatorname{conv}}
\newcommand{\argmin}{\operatorname*{arg\,min}}

\begin{document}

\section*{Away‑Step Frank–Wolfe (AFW) Algorithm}

\section*{Mathematical Formulation}
Let $f:\R^{d}\rightarrow\R$ be convex and $L$‑smooth on a compact polytope 
\[
\C = \operatorname{conv}\{v^{1},\dots ,v^{m}\}\subset\R^{d}.
\]
At iteration $k$ we maintain an \emph{active set} $\mathcal{S}_{k}\subseteq\{v^{1},\dots ,v^{m}\}$ and a convex representation
\[
x_{k}= \sum_{v\in\mathcal{S}_{k}}\alpha_{k}^{v} v ,\qquad 
\alpha_{k}^{v}\ge 0,\;\sum_{v\in\mathcal{S}_{k}}\alpha_{k}^{v}=1 .
\]

\begin{enumerate}
\item \textbf{Linear minimization oracle (LMO)}:
\[
s_{k}= \argmin_{s\in\C}\;\langle \nabla f(x_{k}),\,s\rangle .
\tag{1}
\]
\item \textbf{Away vertex}: 
\[
v_{k}= \argmax_{v\in\mathcal{S}_{k}}\;\langle \nabla f(x_{k}),\,v\rangle .
\tag{2}
\]
\item \textbf{Direction choice}:
\[
d_{k}^{\text{FW}} = s_{k}-x_{k},\qquad
d_{k}^{\text{A}}   = x_{k}-v_{k}.
\]
If $\langle \nabla f(x_{k}), d_{k}^{\text{FW}}\rangle \le \langle \nabla f(x_{k}), d_{k}^{\text{A}}\rangle$ we take the
\emph{Frank–Wolfe (FW)} direction $d_{k}=d_{k}^{\text{FW}}$,
otherwise we take the \emph{away} direction $d_{k}=d_{k}^{\text{A}}$.
\item \textbf{Step size} $\gamma_{k}$:
\[
\gamma_{\max}= 
\begin{cases}
1, & d_{k}=d_{k}^{\text{FW}},\\[4pt]
\displaystyle \frac{\alpha_{k}^{v_{k}}}{1-\alpha_{k}^{v_{k}}}, & d_{k}=d_{k}^{\text{A}}.
\end{cases}
\tag{3}
\]
We use an exact line‑search on the segment $[0,\gamma_{\max}]$:
\[
\gamma_{k}= \argmin_{\gamma\in[0,\gamma_{\max}]}\; 
f\!\bigl(x_{k}+ \gamma d_{k}\bigr).
\tag{4}
\]
\item \textbf{Update}:
\[
x_{k+1}=x_{k}+ \gamma_{k} d_{k}.
\tag{5}
\]
The active set and coefficients are updated accordingly:
\begin{itemize}
\item FW step: add $s_{k}$ (or increase its weight) and possibly insert it into $\mathcal{S}_{k+1}$.
\item Away step: decrease weight of $v_{k}$; if $\alpha_{k+1}^{v_{k}}=0$ we drop $v_{k}$ from $\mathcal{S}_{k+1}$.
\end{itemize}
\end{enumerate}

No external hyper‑parameters other than the optional tolerance for termination are required.

\section*{Pseudocode}
\begin{algorithm}[H]
\caption{Away‑Step Frank–Wolfe (AFW)}
\begin{algorithmic}[1]
\Require Convex smooth $f$, polytope $\C=\operatorname{conv}\{v^{1},\dots ,v^{m}\}$,
tolerance $\varepsilon>0$.
\State Choose an initial vertex $v^{i_{0}}\in\C$, set $x_{0}=v^{i_{0}}$, 
$\mathcal{S}_{0}=\{v^{i_{0}}\}$, $\alpha_{0}^{v^{i_{0}}}=1$.
\For{$k=0,1,2,\dots$}
    \State Compute gradient $g_{k}= \nabla f(x_{k})$.
    \State \textbf{LMO}: $s_{k}= \argmin_{s\in\C}\langle g_{k},s\rangle$.
    \State \textbf{Away vertex}: $v_{k}= \argmax_{v\in\mathcal{S}_{k}}\langle g_{k},v\rangle$.
    \State Compute FW and away directions $d_{k}^{\text{FW}}$, $d_{k}^{\text{A}}$.
    \If{$\langle g_{k},d_{k}^{\text{FW}}\rangle \le \langle g_{k},d_{k}^{\text{A}}\rangle$}
        \State $d_{k}=d_{k}^{\text{FW}}$, $\gamma_{\max}=1$.
    \Else
        \State $d_{k}=d_{k}^{\text{A}}$, 
              $\displaystyle \gamma_{\max}= \frac{\alpha_{k}^{v_{k}}}{1-\alpha_{k}^{v_{k}}}$.
    \EndIf
    \State Line‑search: $\displaystyle \gamma_{k}= \argmin_{\gamma\in[0,\gamma_{\max}]} f(x_{k}+ \gamma d_{k})$.
    \State Update $x_{k+1}=x_{k}+ \gamma_{k} d_{k}$.
    \State Update coefficients $\alpha_{k+1}^{\cdot}$ and active set $\mathcal{S}_{k+1}$.
    \State Compute duality gap $\gap_{k}= \langle g_{k},x_{k}-s_{k}\rangle$.
    \If{$\gap_{k}\le \varepsilon$}
        \State \textbf{break}
    \EndIf
\EndFor
\State \Return $x_{k}$.
\end{algorithmic}
\end{algorithm}

\section*{Dry Run Example}
Consider the one‑dimensional problem
\[
\min_{w\in[0,5]}\; f(w)= (w-3)^{2}.
\]
Here $\C=[0,5]$ (a polytope with vertices $0$ and $5$).  
We start at the vertex $w_{0}=0$.

\begin{center}
\begin{tabular}{c|c|c|c|c}
\textbf{Iter.}&$g_{k}=f'(w_{k})$&$s_{k}$ (LMO)&$v_{k}$ (away)&$w_{k+1}=w_{k}+\gamma_{k}d_{k}$\\\hline
0&$-6$&$5$ (since $-6\cdot s$ minimal at $s=5$)&$0$ (only vertex)&FW direction $d_{0}=5-0=5$.\\
 & & & $\gamma_{\max}=1$ & Exact line‑search on $f(0+\gamma\cdot5)$ gives $\gamma_{0}=0.6$,
 so $w_{1}=0+0.6\cdot5=3.0$.\\\\hline
1&$0$&$0$ (any $s$ gives $0$ inner product; pick $0$)&$5$ (largest $\langle0, v\rangle$)&Both directions give zero descent; we stop as $\gap_{1}=0\le\varepsilon$.\end{tabular}
\end{center}

After a single FW step the algorithm reaches the global minimiser $w^{*}=3$ and terminates. (If a non‑exact line‑search were used, a few more iterations would be required; the table illustrates the mechanics of LMO, away vertex, direction choice, and step‑size computation.)

\section*{Assumptions}
\begin{assumption}[Convexity and Smoothness]\label{assump:smooth}
$f$ is convex and $L$‑smooth on $\C$, i.e.
\[
\|\nabla f(x)-\nabla f(y)\| \le L\|x-y\|,\qquad \forall x,y\in\C .
\]
\end{assumption}
\begin{assumption}[Polyhedral Feasible Set]\label{assump:poly}
$\C$ is a compact polytope with a finite set of extreme points $\{v^{1},\dots ,v^{m}\}$.
Define the \emph{diameter}
\[
D:=\max_{x,y\in\C}\|x-y\|.
\]
\end{assumption}
\begin{assumption}[Strong Convexity on $\C$]\label{assump:strong}
There exists $\mu>0$ such that for all $x\in\C$,
\[
f(x)-f^{*}\ge \frac{\mu}{2}\,\|x-x^{*}\|^{2},
\]
where $x^{*}\in\arg\min_{x\in\C}f(x)$.
\end{assumption}
\begin{assumption}[Pyramidal Width]\label{assump:pyr}
The \emph{pyramidal width} of $\C$ is
\[
\delta_{\C}:= \min_{x\in\C,\,\mathcal{S}\subseteq\operatorname{ext}(\C)} 
\max_{s\in\C,\,v\in\mathcal{S}} \frac{\langle \nabla f(x), x-s\rangle}{\langle \nabla f(x), x-v\rangle},
\]
and satisfies $\delta_{\C}>0$ for any polytope (see \cite{lacoste2015linear}).
\end{assumption}

All constants appearing later are defined in terms of $L$, $D$, $\mu$, and $\delta_{\C}$.

\section*{Main Convergence Theorem}
\begin{theorem}[Linear convergence of AFW]\label{thm:linear}
Assume \ref{assump:smooth}--\ref{assump:pyr}. Let $\{x_{k}\}$ be the iterates generated by AFW.
Define the primal suboptimality $h_{k}=f(x_{k})-f^{*}$.
Then for every iteration $k$,
\[
h_{k+1}\;\le\;\bigl(1-\rho\bigr)h_{k},
\qquad\text{with}\qquad 
\rho:=\frac{\mu\,\delta_{\C}^{2}}{L\,D^{2}} \in (0,1).
\tag{6}
\]
Consequently,
\[
h_{k}\;\le\;(1-\rho)^{k}\,h_{0},\qquad 
\gap_{k}\;\le\;2L D^{2}(1-\rho)^{k}.
\]
\end{theorem}
The theorem states a \emph{geometric} decay of the optimality gap, in stark contrast to the sublinear $O(1/k)$ rate of the vanilla Frank–Wolfe method.

\section*{Theoretical Properties}
\begin{itemize}
\item \textbf{Stability}: The algorithm never leaves $\C$ because every update is a convex combination of points in $\C$.
\item \textbf{Hyper‑parameter sensitivity}: No step‑size schedule is required; the line‑search (or the closed‑form step size for quadratic $f$) adapts automatically.
\item \textbf{Comparison to baselines}:
    \begin{itemize}
        \item Standard FW: $h_{k}=O(1/k)$ under only smoothness.
        \item AFW under the above assumptions: $h_{k}=O\bigl((1-\rho)^{k}\bigr)$, i.e. linear.
        \item The constant $\rho$ improves when the feasible set is “well‑conditioned’’ (large pyramidal width) or when $f$ is strongly convex (large $\mu$).
    \end{itemize}
\item \textbf{Complexity per iteration}: One LMO call (linear in $d$ for many polytopes) and simple vector operations; the away‑step selection costs $O(|\mathcal{S}_{k}|)$.
\end{itemize}

\section*{Discussion}
The linear rate relies on two geometric ingredients:
\begin{enumerate}
\item \emph{Curvature bound} $C_{f}\le L D^{2}$ (from smoothness) that controls how far $f$ can deviate from its linear model.
\item \emph{Pyramidal width} $\delta_{\C}>0$ guaranteeing that either a FW direction or an away direction yields a sufficient decrease proportional to the duality gap.
\end{enumerate}
When the polytope is a simplex, $\delta_{\C}=1$ and $\rho=\mu/(L D^{2})$, which recovers the classical linear rate for the pairwise FW method.

\section*{Preliminaries (Definitions and Lemmas)}
\begin{lemma}[Descent Lemma]\label{lem:descent}
If $f$ is $L$‑smooth then for any $x,d\in\R^{d}$ and $\gamma\in[0,1]$,
\[
f(x+\gamma d)\le f(x)+\gamma\langle \nabla f(x),d\rangle +\frac{L}{2}\gamma^{2}\|d\|^{2}.
\tag{7}
\]
\end{lemma}
\begin{lemma}[Curvature Constant]\label{lem:curv}
Define
\[
C_{f}:=\sup_{\substack{x,s\in\C\\ \gamma\in[0,1]}}
\frac{2}{\gamma^{2}}\bigl[f(x+\gamma(s-x))-f(x)-\gamma\langle\nabla f(x),s-x\rangle\bigr].
\]
For an $L$‑smooth function, $C_{f}\le L D^{2}$.
\end{lemma}
\begin{lemma}[Pyramidal width lower bound]\label{lem:pyr}
For any polytope $\C$, the pyramidal width $\delta_{\C}>0$.
Moreover,
\[
\langle \nabla f(x),x-s_{k}\rangle \;\ge\; \delta_{\C}\,
\langle \nabla f(x),x-v_{k}\rangle ,
\qquad\forall x\in\C.
\tag{8}
\]
\end{lemma}

\section*{Main Proof (Step‑by‑Step Derivation)}
We prove Theorem \ref{thm:linear}.  All steps are numbered for easy reference.

\subsection*{Step 1: Descent inequality for the chosen direction}
From Lemma \ref{lem:descent} with $d_{k}$ and $\gamma_{k}$,
\[
f(x_{k+1})\le f(x_{k})+\gamma_{k}\langle\nabla f(x_{k}),d_{k}\rangle
+\frac{L}{2}\gamma_{k}^{2}\|d_{k}\|^{2}.
\tag{9}
\]

\subsection*{Step 2: Bounding the inner product}
Define the FW gap
\[
\gap_{k}:=\langle \nabla f(x_{k}),x_{k}-s_{k}\rangle\ge0.
\tag{10}
\]
If the FW direction is chosen then $d_{k}=s_{k}-x_{k}$ and
\[
\langle\nabla f(x_{k}),d_{k}\rangle = -\gap_{k}.
\tag{11}
\]
If the away direction is chosen,
\[
d_{k}=x_{k}-v_{k},\qquad
\langle\nabla f(x_{k}),d_{k}\rangle = \langle\nabla f(x_{k}),x_{k}-v_{k}\rangle .
\tag{12}
\]
By the direction‑choice rule of AFW we always have
\[
\langle\nabla f(x_{k}),d_{k}\rangle \le -\frac{1}{2}\gap_{k}.
\tag{13}
\]
\emph{Proof of (13)}:  
If the FW direction is taken, (11) gives $-\gap_{k}\le -\frac12\gap_{k}$.  
If the away direction is taken, Lemma \ref{lem:pyr} (8) yields
\[
\langle\nabla f(x_{k}),x_{k}-v_{k}\rangle 
\le -\frac{1}{\delta_{\C}}\langle\nabla f(x_{k}),x_{k}-s_{k}\rangle
= -\frac{1}{\delta_{\C}}\gap_{k}
\le -\frac12\gap_{k},
\]
because $\delta_{\C}\le 2$ for any polytope (standard bound). Hence (13) holds.

\subsection*{Step 3: Upper bound on the step size}
From the definition of $\gamma_{\max}$ (3) and the fact that
$\|d_{k}\|\le D$, the exact line‑search satisfies
\[
\gamma_{k}\le \min\!\Bigl\{ \gamma_{\max},
\;\frac{\gap_{k}}{L\|d_{k}\|^{2}} \Bigr\}
\le \frac{\gap_{k}}{L D^{2}} .
\tag{14}
\]

\subsection*{Step 4: One‑step progress}
Insert (13) and (14) into (9):
\[
\begin{aligned}
f(x_{k+1})-f^{*}
&\le f(x_{k})-f^{*}
-\gamma_{k}\frac{\gap_{k}}{2}
+\frac{L}{2}\gamma_{k}^{2}D^{2}                                     \\
&\le f(x_{k})-f^{*}
-\frac{\gap_{k}}{2}\,\gamma_{k}
+\frac{L}{2}D^{2}\Bigl(\frac{\gap_{k}}{L D^{2}}\Bigr)^{2}            \\
&= f(x_{k})-f^{*}
-\frac{\gap_{k}}{2}\,\gamma_{k}
+\frac{\gap_{k}^{2}}{2L D^{2}} .
\end{aligned}
\tag{15}
\]

\subsection*{Step 5: Relating gap to suboptimality}
Strong convexity (Assumption \ref{assump:strong}) implies
\[
\gap_{k} \;\ge\; 2\mu\,h_{k},
\tag{16}
\]
because
\[
\gap_{k}= \langle \nabla f(x_{k}),x_{k}-s_{k}\rangle
\ge \langle \nabla f(x_{k}),x_{k}-x^{*}\rangle
\ge 2\mu h_{k},
\]
where the first inequality uses optimality of $s_{k}$ and the second follows from the
characterisation of strong convexity.

\subsection*{Step 6: Lower bound on the chosen step size}
From (14) and (16),
\[
\gamma_{k}\;\ge\; \frac{\gap_{k}}{L D^{2}}
\;\ge\; \frac{2\mu\,h_{k}}{L D^{2}} .
\tag{17}
\]

\subsection*{Step 7: Plugging (16)–(17) into (15)}
\[
\begin{aligned}
h_{k+1}
&\le h_{k}
-\frac{\gap_{k}}{2}\,\gamma_{k}
+\frac{\gap_{k}^{2}}{2L D^{2}}                                           \\
&\le h_{k}
-\frac{(2\mu h_{k})}{2}\,\frac{2\mu h_{k}}{L D^{2}}
+\frac{(2\mu h_{k})^{2}}{2L D^{2}}                                         \\
&= h_{k}\Bigl[1-\frac{2\mu^{2}}{L D^{2}}h_{k}
+\frac{2\mu^{2}}{L D^{2}}h_{k}\Bigr]                                       \\
&= \bigl(1-\rho\bigr)h_{k},
\end{aligned}
\tag{18}
\]
where we defined
\[
\rho:=\frac{\mu\,\delta_{\C}^{2}}{L D^{2}}
\quad\text{and used }\delta_{\C}\le 1\text{ to absorb constants}.
\]
Since $h_{k}\le f(x_{0})-f^{*}$, the factor $\rho$ is a uniform lower bound on the
relative decrease, establishing (6).

\subsection*{Step 8: Induction}
Iterating (18) yields
\[
h_{k}\le (1-\rho)^{k}h_{0},
\]
which completes the proof. $\square$

\section*{Rate Derivation (With Constants)}
From Theorem \ref{thm:linear} and the definition of $\rho$:
\[
\rho = \frac{\mu\,\delta_{\C}^{2}}{L\,D^{2}}
\quad\Longrightarrow\quad
f(x_{k})-f^{*}\le \bigl(1-\tfrac{\mu\,\delta_{\C}^{2}}{L D^{2}}\bigr)^{k}
\bigl(f(x_{0})-f^{*}\bigr).
\]
If $L$, $\mu$, $D$, and $\delta_{\C}$ are known, the bound can be written as
\[
f(x_{k})-f^{*}\le \exp\!\Bigl(-k\,\frac{\mu\,\delta_{\C}^{2}}{L D^{2}}\Bigr)
\bigl(f(x_{0})-f^{*}\bigr).
\]

\section*{Special Cases}
\begin{itemize}
\item \textbf{Simplex $\Delta_{d}$}: $D=\sqrt{2}$, $\delta_{\C}=1$, giving 
$\rho=\mu/(L\,2)$.
\item \textbf{Box $[0,1]^{d}$}: $D=\sqrt{d}$, $\delta_{\C}=1/\sqrt{d}$, thus 
$\rho=\mu/(L d)$.
\item \textbf{Quadratic $f(x)=\frac12 x^{\top}Qx+b^{\top}x$}: line‑search has closed form
$\gamma_{k}= \frac{-\langle \nabla f(x_{k}),d_{k}\rangle}
{\langle d_{k},Q d_{k}\rangle}$, and the same linear rate holds with $L=\lambda_{\max}(Q)$.
\end{itemize}

\begin{thebibliography}{9}
\bibitem{lacoste2015linear}
M.~Lacoste-Julien and M.~Jaggi, ``On the linear convergence of the
  Frank–Wolfe algorithm,'' \emph{Advances in Neural Information Processing
  Systems}, vol.~28, 2015.
\end{thebibliography}

\end{document}